%==============================================================
%  lecons/algebre/surjectivite_exp.tex
%  Surjectivité de l'exponentielle matricielle
%==============================================================

\lecontitle{Surjectivité de l'exponentielle matricielle}{Algèbre linéaire — Logarithme matriciel et formes de Jordan}

%=============================================================
%  § 1 — ÉNONCÉ ET OBSTRUCTIONS
%=============================================================
\secnum{navyblue}{1}{Énoncé et obstructions}

\begin{tcolorbox}[thmbox]
  \textbf{Théorème.}\enspace\index{exponentielle de matrice!surjectivité}
  \begin{enumerate}
    \item \textbf{Sur $\mathbb{C}$.}\enspace
          $\exp : \mathcal{M}_n(\mathbb{C}) \to \mathrm{GL}_n(\mathbb{C})$
          est \emph{surjective}.
    \item \textbf{Sur $\mathbb{R}$.}\enspace
          $\exp : \mathcal{M}_n(\mathbb{R}) \to \mathrm{GL}_n(\mathbb{R})$
          est \emph{non surjective} dès $n \geq 2$.
  \end{enumerate}
\end{tcolorbox}

\rubriq{Obstruction réelle — déterminant}

\begin{mdframed}[style=cexbar]
  La formule de Jacobi\index{formule de Jacobi} $\det(e^A) = e^{\mathrm{tr}(A)}$
  implique $\det(e^A) > 0$ pour tout $A \in \mathcal{M}_n(\mathbb{R})$.
  Ainsi toute matrice à déterminant strictement négatif
  (par exemple $-I_n$ si $n$ est impair) est \emph{hors de l'image} de $\exp$.
\end{mdframed}

\rubriq{Obstruction réelle subtile — Jordan}

\begin{mdframed}[style=cexbar]
  \index{contre-exemple!surjectivité exponentielle réelle}
  La matrice $M = \begin{pmatrix}-1 & 1 \\ 0 & -1\end{pmatrix}$
  satisfait $\det(M) = 1 > 0$, mais \emph{n'est pas dans l'image} de
  $\exp : \mathcal{M}_2(\mathbb{R}) \to \mathrm{GL}_2(\mathbb{R})$.

  \smallskip
  \textit{Preuve.} Supposons $M = e^A$ pour un certain $A \in \mathcal{M}_2(\mathbb{R})$.
  Les valeurs propres de $A$ doivent être des logarithmes de $-1$
  (unique valeur propre de $M$), c'est-à-dire de la forme $(2k+1)\pi i$
  pour $k \in \mathbb{Z}$. Comme $A$ est réelle, ses valeurs propres complexes
  forment des paires conjuguées~: elles seraient $(2k+1)\pi i$ et $-(2k+1)\pi i$,
  deux valeurs distinctes. Donc $A$ est diagonalisable sur $\mathbb{C}$ et
  $e^A = P\,\mathrm{diag}(-1,-1)\,P^{-1} = -I_2 \neq M$. Contradiction.
  \hfill$\blacksquare$
\end{mdframed}

\decsep{}

%=============================================================
%  § 2 — LOGARITHME D'UN BLOC DE JORDAN
%=============================================================
\secnum{navyblue}{2}{Logarithme d'un bloc de Jordan}

\begin{tcolorbox}[thmbox]
  \textbf{Définition (logarithme matriciel).}\enspace%
  \index{logarithme matriciel}
  Pour $\lambda \in \mathbb{C}^*$ et $\mu \in \mathbb{C}$ tel que $e^\mu = \lambda$,
  le \emph{logarithme du bloc de Jordan} $J_m(\lambda)$ est
  \[
    \log J_m(\lambda)
    \;=\;
    \mu\,I_m \;+\; \sum_{k=1}^{m-1} \frac{{(-1)}^{k+1}}{k}\,
    \left(\frac{N}{\lambda}\right)^{\!k},
  \]
  où $N = J_m(\lambda) - \lambda I_m$ est la partie nilpotente ($N^m = 0$).
  La somme est finie.

  \medskip
  \textbf{Propriété.}\enspace $\exp\!\bigl(\log J_m(\lambda)\bigr) = J_m(\lambda)$.
\end{tcolorbox}

\rubriq{Preuve}

\begin{mdframed}[style=proofbar]
  Posons $L_0 = \mu I_m$ et $L_1 = \log(I_m + N/\lambda)
  = \sum_{k=1}^{m-1}\frac{{(-1)}^{k+1}}{k}{(N/\lambda)}^k$ (série formelle
  tronquée car ${(N/\lambda)}^m = 0$).

  $L_0$ et $L_1$ commutent (tous deux polynômes en $N$), donc
  \[
    e^{L_0 + L_1} = e^{L_0}\,e^{L_1} = e^\mu\,\exp\!\bigl(\log(I_m + N/\lambda)\bigr).
  \]
  Or la composition formelle $\exp \circ \log$ donne l'identité sur les
  séries entières autour de $0$~: pour toute matrice nilpotente $Q$,
  $\exp(\log(I+Q)) = I + Q$. Appliqué à $Q = N/\lambda$~:
  \[
    e^{L_0+L_1} = e^\mu\,(I_m + N/\lambda) = \lambda I_m + N = J_m(\lambda).
    \hfill\blacksquare
  \]
\end{mdframed}

\decsep{}

%=============================================================
%  § 3 — PREUVE DE LA SURJECTIVITÉ SUR ℂ
%=============================================================
\secnum{navyblue}{3}{Surjectivité sur $\mathbb{C}$}

\begin{tcolorbox}[thmbox]
  \textbf{Théorème.}\enspace
  Pour tout $M \in \mathrm{GL}_n(\mathbb{C})$, il existe
  $A \in \mathcal{M}_n(\mathbb{C})$ tel que $e^A = M$.
\end{tcolorbox}

\rubriq{Preuve}

\begin{mdframed}[style=proofbar]
  \textbf{Étape 1 — Forme de Jordan.}\enspace
  Il existe $P \in \mathrm{GL}_n(\mathbb{C})$ telle que
  \[
    M = P\,J\,P^{-1}, \qquad
    J = \mathrm{diag}\!\bigl(J_{m_1}(\lambda_1),\ldots,J_{m_r}(\lambda_r)\bigr).
  \]
  Puisque $M$ est inversible, $\lambda_i \neq 0$ pour tout $i$.

  \medskip
  \textbf{Étape 2 — Logarithme bloc par bloc.}\enspace
  Pour chaque $i$, choisir $\mu_i \in \mathbb{C}$ tel que $e^{\mu_i} = \lambda_i$
  et poser $\Lambda_i = \log J_{m_i}(\lambda_i)$ (défini au §\,2). Définir
  \[
    B = \mathrm{diag}(\Lambda_1, \ldots, \Lambda_r), \qquad
    A = P\,B\,P^{-1}.
  \]

  \textbf{Étape 3 — Vérification.}\enspace
  Par conjugaison, $e^A = P\,e^B\,P^{-1}$. De plus,
  \[
    e^B = \mathrm{diag}(e^{\Lambda_1},\ldots,e^{\Lambda_r})
        = \mathrm{diag}(J_{m_1}(\lambda_1),\ldots,J_{m_r}(\lambda_r))
        = J,
  \]
  donc $e^A = PJP^{-1} = M$. \hfill$\blacksquare$
\end{mdframed}

\decsep{}

%=============================================================
%  § 4 — UNICITÉ ET STRUCTURE DU LOGARITHME
%=============================================================
\secnum{navyblue}{4}{Unicité et structure du logarithme}

\begin{tcolorbox}[thmbox]
  \textbf{Proposition.}\enspace
  \begin{enumerate}
    \item Le logarithme matriciel n'est \emph{pas unique}~: les choix de
          $\mu_i$ (branches distinctes de $\log\lambda_i$) donnent
          des antécédents distincts pour un même $M$.
    \item $\exp$ n'est pas injective~: $e^A = e^{A + 2\pi i\,I_n}$,
          et plus généralement $e^A = e^B$ dès que $A$ et $B$ ont
          les mêmes valeurs propres à des multiples de $2\pi i$ près
          \emph{et} le même comportement de Jordan.
    \item L'image de $\exp$ sur $\mathcal{M}_n(\mathbb{R})$ est contenue
          dans $\mathrm{GL}_n^+(\mathbb{R}) = \{M \in \mathrm{GL}_n(\mathbb{R})
          \mid \det M > 0\}$, mais cette inclusion est stricte.
  \end{enumerate}
\end{tcolorbox}

\rubriq{Structure de l'image réelle}

\begin{mdframed}[style=proofbar]
  On a $\det(e^A) = e^{\mathrm{tr}(A)} > 0$, donc
  $\mathrm{Im}(\exp|_{\mathbb{R}}) \subseteq \mathrm{GL}_n^+(\mathbb{R})$.
  L'inclusion est stricte~: le contre-exemple du §\,1 montre que
  $\begin{pmatrix}-1&1\\0&-1\end{pmatrix} \in \mathrm{GL}_2^+(\mathbb{R})$
  n'est pas dans l'image. \hfill$\blacksquare$
\end{mdframed}

\decsep{}

%=============================================================
%  § 5 — PIÈGES ET EXERCICES
%=============================================================
\secnum{exgreen}{5}{Pièges et exercices}

\noindent{\bfseries\color{counterred} Pièges.}

\begin{mdframed}[style=cexbar]
  \textbf{Confondre surjectivité et bijectivité.}\enspace
  Sur $\mathbb{C}$, $\exp$ est surjective mais \emph{pas injective}~:
  le noyau de $\exp$ (en tant qu'homomorphisme de groupes additifs) est
  $2\pi i\,M_n^{\mathrm{diag}}(\mathbb{Z})$, ensemble infini.

  \smallskip\noindent
  \textbf{$\det M > 0$ ne suffit pas sur $\mathbb{R}$.}\enspace
  La condition $\det(e^A)>0$ est nécessaire mais non suffisante pour
  $M \in \mathrm{Im}(\exp|_{\mathbb{R}})$. Le contre-exemple
  $\bigl(\begin{smallmatrix}-1&1\\0&-1\end{smallmatrix}\bigr)$
  est dans $\mathrm{GL}_2^+(\mathbb{R})$ mais hors de l'image.

  \smallskip\noindent
  \textbf{Produit de deux exponentielles.}\enspace
  Tout $M \in \mathrm{GL}_n^+(\mathbb{R})$ est produit \emph{fini}
  d'exponentielles réelles (le groupe de Lie $\mathrm{GL}_n^+(\mathbb{R})$
  est connexe), mais pas nécessairement une seule.
\end{mdframed}

\vspace{6pt}
\exolabel

\begin{mdframed}[style=exobar]
  \begin{enumerate}
    \item Montrer que $-I_n$ est dans l'image de $\exp: \mathcal{M}_n(\mathbb{R})
          \to \mathrm{GL}_n(\mathbb{R})$ si et seulement si $n$ est pair.
          (\textit{Hint~:} pour $n$ pair, utiliser des blocs de rotation~;
          pour $n$ impair, invoquer $\det(e^A) > 0$.)

    \item Calculer explicitement $A \in \mathcal{M}_2(\mathbb{C})$ tel que
          $e^A = \begin{pmatrix}-1&1\\0&-1\end{pmatrix}$.
          (\textit{Hint~:} poser $A = \pi i\,I + B$ et chercher $B$
          nilpotent avec $e^B = -\bigl(\begin{smallmatrix}1&-1\\0&1\end{smallmatrix}\bigr)$.)

    \item Soit $M \in \mathrm{GL}_n(\mathbb{C})$ telle que tous ses blocs de Jordan
          aient des valeurs propres de module $1$. Montrer que $M$ est dans
          l'image de $\exp : \mathcal{M}_n(\mathbb{R})$ si et seulement si
          ses blocs Jordan complexes sont échangés par conjugaison.

    \item Montrer que pour tout $M \in \mathrm{GL}_n(\mathbb{R})$ il existe
          $k \geq 1$ et $A_1, \ldots, A_k \in \mathcal{M}_n(\mathbb{R})$
          tels que $M = e^{A_1}\cdots e^{A_k}$.
          (\textit{Hint~:} décomposer $M = M_+{(-I)}^s$ avec $M_+ \in
          \mathrm{GL}_n^+(\mathbb{R})$ et utiliser $-I_2 = e^{\pi J}$,
          $J = \bigl(\begin{smallmatrix}0&-1\\1&0\end{smallmatrix}\bigr)$.)

    \item (\textit{Défi}) Montrer que
          $\bigl\{e^A \mid A \in \mathcal{M}_2(\mathbb{R})\bigr\} \subsetneq
          \mathrm{GL}_2^+(\mathbb{R})$
          en prouvant que $\bigl(\begin{smallmatrix}-1&1\\0&-1\end{smallmatrix}\bigr)$
          n'est pas dans l'image sans invoquer la forme de Jordan~:
          supposer $e^{tA}$ atteint $M$ en $t=1$ et analyser la trajectoire
          $t \mapsto e^{tA}$ dans $\mathrm{GL}_2(\mathbb{R})$.
  \end{enumerate}
\end{mdframed}

\leconfooter{%
  L.\ Euler (logarithme complexe) —
  É.\ Cartan (groupes de Lie, 1930) —
  R.\ Godement, \textit{Introduction à la théorie des groupes de Lie} (1982)%
}
